% =========================================================
% Limsup / Liminf / Oscillation — Curated Exercise Set
% Sources: Abbott, Bartle & Sherbert, Bruckner, Tao, Lebl
% =========================================================

\section*{Limsup, Liminf, and Oscillation — Exercise Review Set}

The concepts of $\limsup$ and $\liminf$ organize the behavior of arbitrary
sequences. They capture the maximal and minimal subsequential limits and
allow us to analyze oscillatory sequences precisely.

\bigskip

% ---------------------------------------------------------
% Stub Macro
% ---------------------------------------------------------

\newcommand{\exstub}[4]{
\subsection*{Exercise: #1}

\textbf{Source:} #2

\textbf{Technique:} #3

\textbf{Task.} #4

\begin{remark*}[Work Area]
Write your proof or computation here.
\end{remark*}

\bigskip
}

% =========================================================
% ABBOTT
% =========================================================

\section*{Exercises from Abbott}

\exstub
{Abbott 2.5.1}
{Abbott, \textit{Understanding Analysis}}
{unpack-definition}
{Define $\limsup a_n$ and $\liminf a_n$ using tail suprema and infima.}

\exstub
{Abbott 2.5.2}
{Abbott}
{tail supremum argument}
{Show that the sequence $s_n=\sup\{a_k:k\ge n\}$ is decreasing.}

\exstub
{Abbott 2.5.3}
{Abbott}
{tail infimum argument}
{Show that $i_n=\inf\{a_k:k\ge n\}$ is increasing.}

\exstub
{Abbott 2.5.4}
{Abbott}
{limsup construction}
{Prove that $\limsup a_n = \lim s_n$.}

\exstub
{Abbott 2.5.5}
{Abbott}
{liminf construction}
{Prove that $\liminf a_n = \lim i_n$.}

\exstub
{Abbott 2.5.6}
{Abbott}
{subsequence extraction}
{Show that $\limsup a_n$ is the largest subsequential limit.}

\exstub
{Abbott 2.5.7}
{Abbott}
{subsequence extraction}
{Show that $\liminf a_n$ is the smallest subsequential limit.}

% =========================================================
% BARTLE
% =========================================================

\section*{Exercises from Bartle \& Sherbert}

\exstub
{Bartle 3.6.1}
{Bartle \& Sherbert}
{unpack-definition}
{Define limit superior and limit inferior of a sequence.}

\exstub
{Bartle 3.6.2}
{Bartle \& Sherbert}
{tail supremum argument}
{Prove that $\liminf a_n \le \limsup a_n$.}

\exstub
{Bartle 3.6.3}
{Bartle \& Sherbert}
{subsequence extraction}
{Show that there exists a subsequence converging to $\limsup a_n$.}

\exstub
{Bartle 3.6.4}
{Bartle \& Sherbert}
{subsequence extraction}
{Show that there exists a subsequence converging to $\liminf a_n$.}

\exstub
{Bartle 3.6.5}
{Bartle \& Sherbert}
{oscillation bound}
{Show that a sequence converges iff $\limsup a_n = \liminf a_n$.}

\exstub
{Bartle 3.6.6}
{Bartle \& Sherbert}
{oscillation bound}
{Define the oscillation of a sequence and relate it to
$\limsup - \liminf$.}

% =========================================================
% BRUCKNER
% =========================================================

\section*{Exercises from Bruckner}

\exstub
{Bruckner 2.4.1}
{Bruckner, \textit{Real Analysis}}
{unpack-definition}
{Define $\limsup$ and $\liminf$ using subsequential limits.}

\exstub
{Bruckner 2.4.2}
{Bruckner}
{subsequence extraction}
{Show that every subsequential limit lies between
$\liminf a_n$ and $\limsup a_n$.}

\exstub
{Bruckner 2.4.3}
{Bruckner}
{limsup construction}
{Construct a subsequence converging to $\limsup a_n$.}

\exstub
{Bruckner 2.4.4}
{Bruckner}
{liminf construction}
{Construct a subsequence converging to $\liminf a_n$.}

\exstub
{Bruckner 2.4.5}
{Bruckner}
{oscillation bound}
{Show that the oscillation of a convergent sequence is zero.}

\exstub
{Bruckner 2.4.6}
{Bruckner}
{tail supremum argument}
{Prove $\limsup (a_n+b_n) \le \limsup a_n + \limsup b_n$.}

% =========================================================
% TAO
% =========================================================

\section*{Exercises from Tao}

\exstub
{Tao 6.6.1}
{Tao, \textit{Analysis I}}
{unpack-definition}
{Define $\limsup$ using tail suprema.}

\exstub
{Tao 6.6.2}
{Tao}
{subsequence extraction}
{Show $\limsup a_n$ equals the supremum of subsequential limits.}

\exstub
{Tao 6.6.3}
{Tao}
{liminf construction}
{Show $\liminf a_n$ equals the infimum of subsequential limits.}

\exstub
{Tao 6.6.4}
{Tao}
{squeeze argument}
{If $a_n \le b_n \le c_n$ and
$\lim a_n = \lim c_n$, show $\lim b_n$ exists.}

\exstub
{Tao 6.6.5}
{Tao}
{oscillation bound}
{Define the oscillation of a sequence and show convergence
occurs when oscillation tends to zero.}

\exstub
{Tao 6.6.6}
{Tao}
{limsup algebra}
{Prove $\limsup (c a_n)=c\limsup a_n$ for $c>0$.}

% =========================================================
% LEBL
% =========================================================

\section*{Exercises from Lebl}

\exstub
{Lebl 2.4.1}
{Lebl, \textit{Basic Analysis I}}
{unpack-definition}
{Define limit superior and limit inferior of a sequence.}

\exstub
{Lebl 2.4.2}
{Lebl}
{tail supremum argument}
{Show that the sequence of tail suprema is decreasing.}

\exstub
{Lebl 2.4.3}
{Lebl}
{tail infimum argument}
{Show that the sequence of tail infima is increasing.}

\exstub
{Lebl 2.4.4}
{Lebl}
{subsequence extraction}
{Construct a subsequence converging to $\limsup a_n$.}

\exstub
{Lebl 2.4.5}
{Lebl}
{oscillation bound}
{Show that convergence occurs iff $\limsup a_n = \liminf a_n$.}

\exstub
{Lebl 2.4.6}
{Lebl}
{limsup algebra}
{Show $\limsup(a_n+b_n)\le\limsup a_n+\limsup b_n$.}
